\section{Feature dictionary}
\label{sec:appendix:features}

Table~\ref{tab:features} lists every feature in the outfield set, its source, and its
definition. Counting statistics are expressed per ninety minutes played. The two rate
features are used as published rather than divided a second time. Each feature was
confirmed to hold non-null values in both seasons before being admitted, which is the
check that distinguishes a live column from one whose values were withdrawn.

\begin{table}[t]
\centering
\small
\caption{The outfield feature set. Counting statistics are per ninety minutes. The two rate features are used as published.}
\label{tab:features}
\begin{tabular}{llp{7.2cm}}
\toprule
Feature & Source & Definition \\
\midrule
np\_xg\_p90 & Understat & Non-penalty expected goals \\
xa\_p90 & Understat & Expected assists \\
key\_passes\_p90 & Understat & Passes leading directly to a shot \\
xg\_chain\_p90 & Understat & Expected goals of possessions the player was involved in \\
xg\_buildup\_p90 & Understat & As xGChain, excluding shots and key passes \\
shots\_p90 & FBref & Shots taken, excluding penalties \\
shots\_on\_target\_p90 & FBref & Shots on target \\
goals\_non\_penalty\_p90 & FBref & Goals excluding penalties \\
assists\_p90 & FBref & Assists \\
crosses\_p90 & FBref & Crosses attempted \\
interceptions\_padj\_p90 & FBref & Interceptions, possession adjusted \\
tackles\_won\_padj\_p90 & FBref & Tackles won, possession adjusted \\
fouls\_committed\_padj\_p90 & FBref & Fouls committed, possession adjusted \\
fouls\_drawn\_p90 & FBref & Fouls drawn \\
offsides\_p90 & FBref & Times caught offside \\
cards\_yellow\_p90 & FBref & Yellow cards \\
shot\_accuracy\_pct & FBref & Shots on target as a share of shots, a rate \\
goals\_per\_shot & FBref & Goals per shot, a rate \\
\bottomrule
\end{tabular}
\end{table}


\section{Minutes threshold sensitivity}
\label{sec:appendix:sensitivity}

The main analysis excludes outfield players below \MinMinutes{} minutes. That threshold
balances two errors: a low threshold admits players whose per-90 rates are dominated by
sampling noise, while a high threshold discards genuine squad roles such as the impact
substitute. Table~\ref{tab:sensitivity} repeats the entire selection and stability procedure at 270,
\MinMinutes{} and 900 minutes, treating each threshold as an independent pool in which
imputation and standardisation are recomputed and the declared selection rule is reapplied
from scratch. Figure~\ref{fig:sensitivity} plots the same quantities.

The threshold turns out to matter very little. The rule selects two clusters at every
threshold, the silhouette moves only from \SensSilhouetteLow{} to
\SensSilhouetteHigh{}, and bootstrap stability is between \SensARILow{} and
\SensARIHigh{} throughout. More directly, the partition itself barely moves: measured on
the players common to both pools, the adjusted Rand index against the partition at the
threshold actually used is \SensAgreeLow{} at 270 minutes and \SensAgreeHigh{} at 900. A
reader who preferred a different threshold would reach the same conclusions.

\begin{table}[t]
\centering
\small
\caption{Sensitivity of the analysis to the minutes threshold in the primary season.}
\label{tab:sensitivity}
\begin{tabular}{rrrrr}
\toprule
Minutes & Eligible & $k$ & Silhouette & Bootstrap ARI \\
\midrule
270 & 403 & n/a & n/a & n/a \\
450 & 366 & n/a & n/a & n/a \\
900 & 296 & n/a & n/a & n/a \\
\bottomrule
\end{tabular}
\end{table}


\begin{figure}[h]
  \centering
  \includegraphics[width=\textwidth]{minutes_sensitivity.pdf}
  \caption{Eligible sample, silhouette and bootstrap stability across the candidate
  minutes thresholds. The dashed line marks the threshold used in the main analysis.}
  \label{fig:sensitivity}
\end{figure}

\section{UMAP hyperparameter grid}
\label{sec:appendix:umap}

Table~\ref{tab:umapgrid} reports trustworthiness for all twelve configurations of the
neighbour count and minimum distance. Selection was by trustworthiness alone, decided
before any embedding was inspected visually.

\begin{table}[t]
\centering
\small
\caption{UMAP hyperparameter grid, ranked by trustworthiness. Selection was by trustworthiness alone and was decided before any embedding was inspected.}
\label{tab:umapgrid}
\begin{tabular}{rrr}
\toprule
Neighbours & Minimum distance & Trustworthiness \\
\midrule
10 & 0 & 0.9337 \\
15 & 0 & 0.9322 \\
15 & 0.10 & 0.9318 \\
30 & 0 & 0.9312 \\
10 & 0.10 & 0.9309 \\
50 & 0 & 0.9273 \\
30 & 0.10 & 0.9242 \\
10 & 0.50 & 0.9218 \\
50 & 0.10 & 0.9179 \\
15 & 0.50 & 0.9134 \\
50 & 0.50 & 0.9091 \\
30 & 0.50 & 0.9084 \\
\bottomrule
\end{tabular}
\end{table}


\section{Cluster selection criteria in full}
\label{sec:appendix:kselection}

Table~\ref{tab:kselection} gives the silhouette coefficient, Davies-Bouldin index,
Calinski-Harabasz index and gap statistic for every candidate number of clusters, for the
global analysis and within each position group. The consensus rule described in
Section~\ref{sec:methods} is applied to these values.

\begin{longtable}{llrrrrr}
\caption{Cluster selection criteria for every candidate number of clusters, by scope. The consensus rule is applied to these values.}\\
\label{tab:kselection}\\
\toprule
Scope & $k$ & Silhouette & Davies-Bouldin & Calinski-Harabasz & Gap \\
\midrule
\endfirsthead
\toprule
Scope & $k$ & Silhouette & Davies-Bouldin & Calinski-Harabasz & Gap \\
\midrule
\endhead
global$\ast$ & 2 & 0.287 & 1.451 & 145.6 & 1.216 \\
global & 3 & 0.215 & 1.695 & 101.9 & 1.283 \\
global & 4 & 0.154 & 1.954 & 86.1 & 1.339 \\
global & 5 & 0.145 & 1.861 & 75.4 & 1.381 \\
global & 6 & 0.144 & 1.797 & 68.3 & 1.415 \\
global & 7 & 0.135 & 1.905 & 61.3 & 1.438 \\
global & 8 & 0.118 & 1.881 & 56.1 & 1.458 \\
global & 9 & 0.123 & 1.840 & 53.5 & 1.477 \\
global & 10 & 0.123 & 1.805 & 49.8 & 1.484 \\
global & 11 & 0.128 & 1.657 & 47.6 & 1.506 \\
global & 12 & 0.121 & 1.911 & 44.6 & 1.517 \\
DF$\ast$ & 2 & 0.204 & 2.037 & 25.8 & 0.943 \\
DF & 3 & 0.151 & 1.919 & 23.5 & 1.011 \\
DF & 4 & 0.142 & 1.983 & 19.7 & 1.036 \\
DF & 5 & 0.120 & 1.961 & 18.0 & 1.059 \\
DF & 6 & 0.117 & 1.795 & 16.9 & 1.083 \\
DF & 7 & 0.106 & 1.898 & 15.6 & 1.109 \\
DF & 8 & 0.120 & 1.792 & 14.7 & 1.136 \\
DF & 9 & 0.120 & 1.754 & 14.1 & 1.154 \\
DF & 10 & 0.112 & 1.636 & 13.5 & 1.159 \\
DF & 11 & 0.098 & 1.636 & 12.8 & 1.162 \\
DF & 12 & 0.105 & 1.627 & 12.3 & 1.196 \\
MF$\ast$ & 2 & 0.243 & 1.543 & 69.4 & 1.035 \\
MF & 3 & 0.147 & 1.969 & 49.0 & 1.091 \\
MF & 4 & 0.145 & 1.895 & 40.4 & 1.130 \\
MF & 5 & 0.145 & 1.873 & 34.7 & 1.153 \\
MF & 6 & 0.132 & 1.939 & 30.3 & 1.167 \\
MF & 7 & 0.122 & 2.054 & 27.4 & 1.199 \\
MF & 8 & 0.107 & 2.030 & 24.9 & 1.202 \\
MF & 9 & 0.103 & 1.923 & 23.8 & 1.212 \\
MF & 10 & 0.106 & 1.983 & 21.9 & 1.220 \\
MF & 11 & 0.094 & 1.944 & 20.4 & 1.227 \\
MF & 12 & 0.097 & 1.835 & 19.5 & 1.226 \\
FW$\ast$ & 2 & 0.187 & 1.803 & 12.4 & 0.764 \\
FW & 3 & 0.175 & 1.576 & 10.8 & 0.837 \\
FW & 4 & 0.176 & 1.599 & 10.2 & 0.908 \\
FW & 5 & 0.196 & 1.181 & 9.3 & 0.909 \\
FW & 6 & 0.124 & 1.490 & 8.5 & 0.944 \\
FW & 7 & 0.132 & 1.350 & 7.8 & 1.017 \\
FW & 8 & 0.101 & 1.389 & 7.7 & 1.021 \\
FW & 9 & 0.127 & 1.229 & 7.7 & 1.070 \\
FW & 10 & 0.125 & 1.350 & 7.8 & 1.051 \\
FW & 11 & 0.128 & 1.199 & 7.1 & 1.064 \\
FW & 12 & 0.116 & 1.126 & 6.8 & 1.096 \\
\bottomrule
\end{longtable}


\section{Cluster membership}
\label{sec:appendix:membership}

Table~\ref{tab:membership} lists every eligible outfield player in the earlier of the two reduced-sample seasons
with his club, listed position, assigned archetype and distance from his cluster
centroid.

\small
\begin{longtable}{p{3.4cm}p{2.5cm}p{1.0cm}p{5.0cm}r}
\caption{Cluster membership for every eligible outfield player in the primary season.}\\
\label{tab:membership}\\
\toprule
Player & Club & Listed & Archetype & Distance \\
\midrule
\endfirsthead
\toprule
Player & Club & Listed & Archetype & Distance \\
\midrule
\endhead
Aaron Wan-Bissaka & West Ham & DF & Crossing and chance creating defenders & 3.09 \\
Andy Robertson & Liverpool & DF & Crossing and chance creating defenders & 4.59 \\
Antonee Robinson & Fulham & DF & Crossing and chance creating defenders & 4.28 \\
Ashley Young & Everton & DF & Crossing and chance creating defenders & 4.02 \\
Ben White & Arsenal & DF & Crossing and chance creating defenders & 2.68 \\
Conor Bradley & Liverpool & DF & Crossing and chance creating defenders & 6.54 \\
Cristian Romero & Tottenham & DF & Crossing and chance creating defenders & 4.57 \\
Dean Huijsen & Bournemouth & DF & Crossing and chance creating defenders & 3.65 \\
Destiny Udogie & Tottenham & DF & Crossing and chance creating defenders & 4.51 \\
Djed Spence & Tottenham & DF & Crossing and chance creating defenders & 2.51 \\
Ian Maatsen & Aston Villa & DF & Crossing and chance creating defenders & 5.18 \\
Joško Gvardiol & Manchester City & DF & Crossing and chance creating defenders & 4.51 \\
Jurriën Timber & Arsenal & DF & Crossing and chance creating defenders & 3.25 \\
Keane Lewis-Potter & Brentford & DF & Crossing and chance creating defenders & 3.78 \\
Kenny Tete & Fulham & DF & Crossing and chance creating defenders & 4.14 \\
Kieran Trippier & Newcastle & DF & Crossing and chance creating defenders & 4.42 \\
Leif Davis & Ipswich Town & DF & Crossing and chance creating defenders & 4.61 \\
Lewis Hall & Newcastle & DF & Crossing and chance creating defenders & 3.95 \\
Lisandro Martínez & Manchester Utd & DF & Crossing and chance creating defenders & 3.95 \\
Lucas Digne & Aston Villa & DF & Crossing and chance creating defenders & 4.50 \\
Mads Roerslev & Brentford & DF & Crossing and chance creating defenders & 4.53 \\
Malo Gusto & Chelsea & DF & Crossing and chance creating defenders & 2.37 \\
Marc Cucurella & Chelsea & DF & Crossing and chance creating defenders & 3.73 \\
Matheus Nunes & Manchester City & DF & Crossing and chance creating defenders & 5.13 \\
Mats Wieffer & Brighton & DF & Crossing and chance creating defenders & 6.15 \\
Michael Kayode & Brentford & DF & Crossing and chance creating defenders & 4.18 \\
Milos Kerkez & Bournemouth & DF & Crossing and chance creating defenders & 2.82 \\
Neco Williams & Nottingham & DF & Crossing and chance creating defenders & 3.49 \\
Nico O'Reilly & Manchester City & DF & Crossing and chance creating defenders & 7.55 \\
Oleksandr Zinchenko & Arsenal & DF & Crossing and chance creating defenders & 5.56 \\
Pedro Porro & Tottenham & DF & Crossing and chance creating defenders & 4.60 \\
Pervis Estupiñán & Brighton & DF & Crossing and chance creating defenders & 2.76 \\
Reece James & Chelsea & DF & Crossing and chance creating defenders & 2.78 \\
Riccardo Calafiori & Arsenal & DF & Crossing and chance creating defenders & 5.54 \\
Ryan Manning & Southampton & DF & Crossing and chance creating defenders & 5.20 \\
Tariq Lamptey & Brighton & DF & Crossing and chance creating defenders & 5.15 \\
Trent Alexander-Arnold & Liverpool & DF & Crossing and chance creating defenders & 6.44 \\
Trevoh Chalobah & Crystal Palace / Chelsea & DF & Crossing and chance creating defenders & 3.86 \\
Aaron Cresswell & West Ham & DF & Low attacking involvement defenders & 4.53 \\
Abdukodir Khusanov & Manchester City & DF & Low attacking involvement defenders & 3.38 \\
Adam Smith & Bournemouth & DF & Low attacking involvement defenders & 3.13 \\
Adam Webster & Brighton & DF & Low attacking involvement defenders & 4.13 \\
Archie Gray & Tottenham & DF & Low attacking involvement defenders & 4.44 \\
Axel Disasi & Aston Villa / Chelsea & DF & Low attacking involvement defenders & 3.55 \\
Axel Tuanzebe & Ipswich Town & DF & Low attacking involvement defenders & 2.69 \\
Ben Davies & Tottenham & DF & Low attacking involvement defenders & 2.29 \\
Caleb Okoli & Leicester City & DF & Low attacking involvement defenders & 3.05 \\
Calvin Bassey & Fulham & DF & Low attacking involvement defenders & 1.45 \\
Cameron Burgess & Ipswich Town & DF & Low attacking involvement defenders & 3.22 \\
Chris Richards & Crystal Palace & DF & Low attacking involvement defenders & 2.65 \\
Conor Coady & Leicester City & DF & Low attacking involvement defenders & 5.36 \\
Craig Dawson & Wolves & DF & Low attacking involvement defenders & 2.87 \\
Dan Burn & Newcastle & DF & Low attacking involvement defenders & 2.06 \\
Dara O'Shea & Ipswich Town & DF & Low attacking involvement defenders & 2.10 \\
Diego Carlos & Aston Villa & DF & Low attacking involvement defenders & 3.36 \\
Emerson Palmieri & West Ham & DF & Low attacking involvement defenders & 3.18 \\
Emmanuel Agbadou & Wolves & DF & Low attacking involvement defenders & 1.68 \\
Ethan Pinnock & Brentford & DF & Low attacking involvement defenders & 3.21 \\
Ezri Konsa & Aston Villa & DF & Low attacking involvement defenders & 3.18 \\
Fabian Schär & Newcastle & DF & Low attacking involvement defenders & 4.54 \\
Gabriel Magalhães & Arsenal & DF & Low attacking involvement defenders & 4.17 \\
Harry Maguire & Manchester Utd & DF & Low attacking involvement defenders & 4.80 \\
Ibrahima Konaté & Liverpool & DF & Low attacking involvement defenders & 3.34 \\
Igor & Brighton & DF & Low attacking involvement defenders & 3.80 \\
Illia Zabarnyi & Bournemouth & DF & Low attacking involvement defenders & 2.39 \\
Issa Diop & Fulham & DF & Low attacking involvement defenders & 1.92 \\
Jack Stephens & Southampton & DF & Low attacking involvement defenders & 2.43 \\
Jacob Greaves & Ipswich Town & DF & Low attacking involvement defenders & 1.92 \\
Jake O'Brien & Everton & DF & Low attacking involvement defenders & 4.00 \\
Jakub Kiwior & Arsenal & DF & Low attacking involvement defenders & 3.05 \\
James Bree & Southampton & DF & Low attacking involvement defenders & 3.61 \\
James Justin & Leicester City & DF & Low attacking involvement defenders & 2.82 \\
James Tarkowski & Everton & DF & Low attacking involvement defenders & 2.08 \\
Jan Bednarek & Southampton & DF & Low attacking involvement defenders & 4.44 \\
Jan Paul van Hecke & Brighton & DF & Low attacking involvement defenders & 1.75 \\
Jannik Vestergaard & Leicester City & DF & Low attacking involvement defenders & 2.96 \\
Jarell Quansah & Liverpool & DF & Low attacking involvement defenders & 4.41 \\
Jarrad Branthwaite & Everton & DF & Low attacking involvement defenders & 2.12 \\
Jean-Clair Todibo & West Ham & DF & Low attacking involvement defenders & 2.98 \\
Joachim Andersen & Fulham / Crystal Palace & DF & Low attacking involvement defenders & 2.00 \\
Joe Gomez & Liverpool & DF & Low attacking involvement defenders & 5.67 \\
John Stones & Manchester City & DF & Low attacking involvement defenders & 9.96 \\
Joël Veltman & Brighton & DF & Low attacking involvement defenders & 4.56 \\
Julian Araujo & Bournemouth & DF & Low attacking involvement defenders & 3.64 \\
Kevin Danso & Tottenham & DF & Low attacking involvement defenders & 3.62 \\
Konstantinos Mavropanos & West Ham & DF & Low attacking involvement defenders & 3.24 \\
Kostas Tsimikas & Liverpool & DF & Low attacking involvement defenders & 6.15 \\
Kristoffer Ajer & Brentford & DF & Low attacking involvement defenders & 2.55 \\
Kyle Walker & Manchester City & DF & Low attacking involvement defenders & 4.22 \\
Lamare Bogarde & Aston Villa & DF & Low attacking involvement defenders & 5.84 \\
Leny Yoro & Manchester Utd & DF & Low attacking involvement defenders & 2.85 \\
Levi Colwill & Chelsea & DF & Low attacking involvement defenders & 2.38 \\
Lewis Dunk & Brighton & DF & Low attacking involvement defenders & 2.63 \\
Luke Thomas & Leicester City & DF & Low attacking involvement defenders & 4.26 \\
Luke Woolfenden & Ipswich Town & DF & Low attacking involvement defenders & 4.15 \\
Manuel Akanji & Manchester City & DF & Low attacking involvement defenders & 3.82 \\
Marc Guéhi & Crystal Palace & DF & Low attacking involvement defenders & 2.44 \\
Marcos Senesi & Bournemouth & DF & Low attacking involvement defenders & 4.07 \\
Matt Doherty & Wolves & DF & Low attacking involvement defenders & 1.89 \\
Matthijs de Ligt & Manchester Utd & DF & Low attacking involvement defenders & 4.15 \\
Matty Cash & Aston Villa & DF & Low attacking involvement defenders & 4.01 \\
Max Kilman & West Ham & DF & Low attacking involvement defenders & 1.82 \\
Maxence Lacroix & Crystal Palace & DF & Low attacking involvement defenders & 2.59 \\
Michael Keane & Everton & DF & Low attacking involvement defenders & 5.83 \\
Micky van de Ven & Tottenham & DF & Low attacking involvement defenders & 3.68 \\
Morato & Nottingham & DF & Low attacking involvement defenders & 3.67 \\
Murillo & Nottingham & DF & Low attacking involvement defenders & 2.20 \\
Myles Lewis-Skelly & Arsenal & DF & Low attacking involvement defenders & 5.97 \\
Nathan Aké & Manchester City & DF & Low attacking involvement defenders & 3.94 \\
Nathan Collins & Brentford & DF & Low attacking involvement defenders & 2.35 \\
Nathan Wood-Gordon & Southampton & DF & Low attacking involvement defenders & 2.74 \\
Nathaniel Clyne & Crystal Palace & DF & Low attacking involvement defenders & 3.16 \\
Nikola Milenković & Nottingham & DF & Low attacking involvement defenders & 3.91 \\
Noussair Mazraoui & Manchester Utd & DF & Low attacking involvement defenders & 3.75 \\
Ola Aina & Nottingham & DF & Low attacking involvement defenders & 2.21 \\
Pau Torres & Aston Villa & DF & Low attacking involvement defenders & 2.70 \\
Radu Drăgușin & Tottenham & DF & Low attacking involvement defenders & 2.49 \\
Rico Lewis & Manchester City & DF & Low attacking involvement defenders & 4.39 \\
Rúben Dias & Manchester City & DF & Low attacking involvement defenders & 3.75 \\
Santiago Bueno & Wolves & DF & Low attacking involvement defenders & 2.97 \\
Sepp van den Berg & Brentford & DF & Low attacking involvement defenders & 2.18 \\
Taylor Harwood-Bellis & Southampton & DF & Low attacking involvement defenders & 2.28 \\
Timothy Castagne & Fulham & DF & Low attacking involvement defenders & 2.93 \\
Tino Livramento & Newcastle & DF & Low attacking involvement defenders & 3.48 \\
Tosin Adarabioyo & Chelsea & DF & Low attacking involvement defenders & 3.27 \\
Toti Gomes & Wolves & DF & Low attacking involvement defenders & 2.24 \\
Tyrone Mings & Aston Villa & DF & Low attacking involvement defenders & 3.08 \\
Victor Bernth Kristiansen & Leicester City & DF & Low attacking involvement defenders & 4.22 \\
Victor Lindelöf & Manchester Utd & DF & Low attacking involvement defenders & 3.45 \\
Virgil van Dijk & Liverpool & DF & Low attacking involvement defenders & 4.08 \\
Vitaliy Mykolenko & Everton & DF & Low attacking involvement defenders & 3.14 \\
Vladimír Coufal & West Ham & DF & Low attacking involvement defenders & 4.23 \\
Wesley Fofana & Chelsea & DF & Low attacking involvement defenders & 4.97 \\
William Saliba & Arsenal & DF & Low attacking involvement defenders & 3.40 \\
Wout Faes & Leicester City & DF & Low attacking involvement defenders & 2.33 \\
Álex Moreno & Nottingham & DF & Low attacking involvement defenders & 3.76 \\
Adam Armstrong & Southampton & FW & Crossing and creating wide forwards & 3.69 \\
Anthony Gordon & Newcastle & FW & Crossing and creating wide forwards & 2.96 \\
Bukayo Saka & Arsenal & FW & Crossing and creating wide forwards & 6.23 \\
Ethan Nwaneri & Arsenal & FW & Crossing and creating wide forwards & 2.98 \\
Fabio Carvalho & Brentford & FW & Crossing and creating wide forwards & 5.32 \\
Gabriel Martinelli & Arsenal & FW & Crossing and creating wide forwards & 2.17 \\
Harvey Barnes & Newcastle & FW & Crossing and creating wide forwards & 2.69 \\
Jacob Murphy & Newcastle & FW & Crossing and creating wide forwards & 3.43 \\
Joshua Zirkzee & Manchester Utd & FW & Crossing and creating wide forwards & 3.12 \\
João Pedro & Brighton & FW & Crossing and creating wide forwards & 3.29 \\
Kamaldeen Sulemana & Southampton & FW & Crossing and creating wide forwards & 4.64 \\
Leandro Trossard & Arsenal & FW & Crossing and creating wide forwards & 1.28 \\
Luis Díaz & Liverpool & FW & Crossing and creating wide forwards & 4.40 \\
Mathys Tel & Tottenham & FW & Crossing and creating wide forwards & 2.45 \\
Michail Antonio & West Ham & FW & Crossing and creating wide forwards & 3.54 \\
Raheem Sterling & Arsenal & FW & Crossing and creating wide forwards & 5.73 \\
Son Heung-min & Tottenham & FW & Crossing and creating wide forwards & 3.42 \\
Timo Werner & Tottenham & FW & Crossing and creating wide forwards & 5.35 \\
Wilson Odobert & Tottenham & FW & Crossing and creating wide forwards & 4.20 \\
Alexander Isak & Newcastle & FW & Shot taking central forwards & 3.60 \\
Beto & Everton & FW & Shot taking central forwards & 2.05 \\
Brennan Johnson & Tottenham & FW & Shot taking central forwards & 2.55 \\
Cameron Archer & Southampton & FW & Shot taking central forwards & 4.03 \\
Chris Wood & Nottingham & FW & Shot taking central forwards & 3.64 \\
Danny Welbeck & Brighton & FW & Shot taking central forwards & 2.85 \\
Darwin Núñez & Liverpool & FW & Shot taking central forwards & 5.20 \\
Diogo Jota & Liverpool & FW & Shot taking central forwards & 5.19 \\
Dominic Calvert-Lewin & Everton & FW & Shot taking central forwards & 4.01 \\
Dominic Solanke & Tottenham & FW & Shot taking central forwards & 2.44 \\
Eddie Nketiah & Crystal Palace & FW & Shot taking central forwards & 4.57 \\
Erling Haaland & Manchester City & FW & Shot taking central forwards & 4.19 \\
Evanilson & Bournemouth & FW & Shot taking central forwards & 2.19 \\
Gabriel Jesus & Arsenal & FW & Shot taking central forwards & 5.68 \\
George Hirst & Ipswich Town & FW & Shot taking central forwards & 2.57 \\
Jamie Vardy & Leicester City & FW & Shot taking central forwards & 3.08 \\
Jean-Philippe Mateta & Crystal Palace & FW & Shot taking central forwards & 2.09 \\
Jáder Durán & Aston Villa & FW & Shot taking central forwards & 7.28 \\
Jørgen Strand Larsen & Wolves & FW & Shot taking central forwards & 3.19 \\
Kai Havertz & Arsenal & FW & Shot taking central forwards & 1.80 \\
Liam Delap & Ipswich Town & FW & Shot taking central forwards & 2.81 \\
Niclas Füllkrug & West Ham & FW & Shot taking central forwards & 4.18 \\
Nicolas Jackson & Chelsea & FW & Shot taking central forwards & 1.92 \\
Ollie Watkins & Aston Villa & FW & Shot taking central forwards & 2.30 \\
Paul Onuachu & Southampton & FW & Shot taking central forwards & 2.75 \\
Rasmus Højlund & Manchester Utd & FW & Shot taking central forwards & 4.07 \\
Raúl Jiménez & Fulham & FW & Shot taking central forwards & 2.60 \\
Richarlison & Tottenham & FW & Shot taking central forwards & 3.79 \\
Rodrigo Muniz & Fulham & FW & Shot taking central forwards & 4.35 \\
Yoane Wissa & Brentford & FW & Shot taking central forwards & 2.27 \\
Abdoulaye Doucouré & Everton & MF & Ball winning, low attacking involvement midfielders & 3.11 \\
Adam Wharton & Crystal Palace & MF & Ball winning, low attacking involvement midfielders & 3.29 \\
Alex Scott & Bournemouth & MF & Ball winning, low attacking involvement midfielders & 4.20 \\
Amadou Onana & Aston Villa & MF & Ball winning, low attacking involvement midfielders & 3.37 \\
André & Wolves & MF & Ball winning, low attacking involvement midfielders & 2.86 \\
Ben Johnson & Ipswich Town & MF & Ball winning, low attacking involvement midfielders & 2.89 \\
Bilal El Khannouss & Leicester City & MF & Ball winning, low attacking involvement midfielders & 3.33 \\
Bobby De Cordova-Reid & Leicester City & MF & Ball winning, low attacking involvement midfielders & 4.65 \\
Boubacar Kamara & Aston Villa & MF & Ball winning, low attacking involvement midfielders & 2.30 \\
Boubakary Soumaré & Leicester City & MF & Ball winning, low attacking involvement midfielders & 3.36 \\
Bruno Guimarães & Newcastle & MF & Ball winning, low attacking involvement midfielders & 3.65 \\
Carlos Baleba & Brighton & MF & Ball winning, low attacking involvement midfielders & 2.81 \\
Carlos Soler & West Ham & MF & Ball winning, low attacking involvement midfielders & 2.53 \\
Casemiro & Manchester Utd & MF & Ball winning, low attacking involvement midfielders & 4.01 \\
Cheick Doucouré & Crystal Palace & MF & Ball winning, low attacking involvement midfielders & 5.22 \\
Christian Nørgaard & Brentford & MF & Ball winning, low attacking involvement midfielders & 2.64 \\
Conor Chaplin & Ipswich Town & MF & Ball winning, low attacking involvement midfielders & 2.92 \\
Daichi Kamada & Crystal Palace & MF & Ball winning, low attacking involvement midfielders & 1.75 \\
Daniel Muñoz & Crystal Palace & MF & Ball winning, low attacking involvement midfielders & 3.51 \\
Diogo Dalot & Manchester Utd & MF & Ball winning, low attacking involvement midfielders & 2.19 \\
Edson Álvarez & West Ham & MF & Ball winning, low attacking involvement midfielders & 3.08 \\
Elliot Anderson & Nottingham & MF & Ball winning, low attacking involvement midfielders & 2.18 \\
Flynn Downes & Southampton & MF & Ball winning, low attacking involvement midfielders & 2.78 \\
Guido Rodríguez & West Ham & MF & Ball winning, low attacking involvement midfielders & 2.80 \\
Harry Winks & Leicester City & MF & Ball winning, low attacking involvement midfielders & 3.17 \\
Hwang Hee-chan & Wolves & MF & Ball winning, low attacking involvement midfielders & 6.60 \\
Ibrahim Sangaré & Nottingham & MF & Ball winning, low attacking involvement midfielders & 6.14 \\
Idrissa Gueye & Everton & MF & Ball winning, low attacking involvement midfielders & 2.34 \\
Jack Clarke & Ipswich Town & MF & Ball winning, low attacking involvement midfielders & 4.03 \\
Jack Harrison & Everton & MF & Ball winning, low attacking involvement midfielders & 4.30 \\
Jack Hinshelwood & Brighton & MF & Ball winning, low attacking involvement midfielders & 4.61 \\
Jack Taylor & Ipswich Town & MF & Ball winning, low attacking involvement midfielders & 2.43 \\
Jacob Ramsey & Aston Villa & MF & Ball winning, low attacking involvement midfielders & 2.84 \\
Jaden Philogene Bidace & Ipswich Town / Aston Villa & MF & Ball winning, low attacking involvement midfielders & 4.28 \\
James Garner & Everton & MF & Ball winning, low attacking involvement midfielders & 2.04 \\
James Ward-Prowse & West Ham / Nottingham & MF & Ball winning, low attacking involvement midfielders & 3.03 \\
Jean-Ricner Bellegarde & Wolves & MF & Ball winning, low attacking involvement midfielders & 4.06 \\
Jefferson Lerma & Crystal Palace & MF & Ball winning, low attacking involvement midfielders & 2.61 \\
Jens Cajuste & Ipswich Town & MF & Ball winning, low attacking involvement midfielders & 2.83 \\
Jesper Lindstrøm & Everton & MF & Ball winning, low attacking involvement midfielders & 3.68 \\
Joe Aribo & Southampton & MF & Ball winning, low attacking involvement midfielders & 3.13 \\
Joe Willock & Newcastle & MF & Ball winning, low attacking involvement midfielders & 3.56 \\
Joelinton & Newcastle & MF & Ball winning, low attacking involvement midfielders & 3.01 \\
John McGinn & Aston Villa & MF & Ball winning, low attacking involvement midfielders & 3.29 \\
Jordan Ayew & Leicester City / Crystal Palace & MF & Ball winning, low attacking involvement midfielders & 4.60 \\
Jorginho & Arsenal & MF & Ball winning, low attacking involvement midfielders & 4.85 \\
João Gomes & Wolves & MF & Ball winning, low attacking involvement midfielders & 2.56 \\
Kalvin Phillips & Ipswich Town & MF & Ball winning, low attacking involvement midfielders & 2.81 \\
Kasey McAteer & Leicester City & MF & Ball winning, low attacking involvement midfielders & 3.50 \\
Kobbie Mainoo & Manchester Utd & MF & Ball winning, low attacking involvement midfielders & 1.93 \\
Kyle Walker-Peters & Southampton & MF & Ball winning, low attacking involvement midfielders & 2.29 \\
Lesley Ugochukwu & Southampton & MF & Ball winning, low attacking involvement midfielders & 3.93 \\
Lewis Cook & Bournemouth & MF & Ball winning, low attacking involvement midfielders & 2.67 \\
Lucas Bergvall & Tottenham & MF & Ball winning, low attacking involvement midfielders & 2.27 \\
Lucas Paquetá & West Ham & MF & Ball winning, low attacking involvement midfielders & 2.85 \\
Manuel Ugarte & Manchester Utd & MF & Ball winning, low attacking involvement midfielders & 3.67 \\
Mario Lemina & Wolves & MF & Ball winning, low attacking involvement midfielders & 3.07 \\
Mateo Kovačić & Manchester City & MF & Ball winning, low attacking involvement midfielders & 3.82 \\
Mateus Fernandes & Southampton & MF & Ball winning, low attacking involvement midfielders & 1.86 \\
Mathias Jensen & Brentford & MF & Ball winning, low attacking involvement midfielders & 3.67 \\
Moisés Caicedo & Chelsea & MF & Ball winning, low attacking involvement midfielders & 3.08 \\
Nicolás Domínguez & Nottingham & MF & Ball winning, low attacking involvement midfielders & 2.76 \\
Nicolás González & Manchester City & MF & Ball winning, low attacking involvement midfielders & 5.07 \\
Nélson Semedo & Wolves & MF & Ball winning, low attacking involvement midfielders & 1.85 \\
Oliver Scarles & West Ham & MF & Ball winning, low attacking involvement midfielders & 4.31 \\
Oliver Skipp & Leicester City & MF & Ball winning, low attacking involvement midfielders & 4.02 \\
Omari Hutchinson & Ipswich Town & MF & Ball winning, low attacking involvement midfielders & 3.30 \\
Orel Mangala & Everton & MF & Ball winning, low attacking involvement midfielders & 2.87 \\
Pape Matar Sarr & Tottenham & MF & Ball winning, low attacking involvement midfielders & 2.51 \\
Patrick Dorgu & Manchester Utd & MF & Ball winning, low attacking involvement midfielders & 4.54 \\
Patson Daka & Leicester City & MF & Ball winning, low attacking involvement midfielders & 4.79 \\
Rayan Aït-Nouri & Wolves & MF & Ball winning, low attacking involvement midfielders & 2.51 \\
Rodrigo Bentancur & Tottenham & MF & Ball winning, low attacking involvement midfielders & 4.07 \\
Rodrigo Gomes & Wolves & MF & Ball winning, low attacking involvement midfielders & 5.18 \\
Roméo Lavia & Chelsea & MF & Ball winning, low attacking involvement midfielders & 4.01 \\
Ryan Christie & Bournemouth & MF & Ball winning, low attacking involvement midfielders & 2.35 \\
Ryan Gravenberch & Liverpool & MF & Ball winning, low attacking involvement midfielders & 3.56 \\
Ryan Yates & Nottingham & MF & Ball winning, low attacking involvement midfielders & 3.48 \\
Sam Morsy & Ipswich Town & MF & Ball winning, low attacking involvement midfielders & 1.72 \\
Sander Berge & Fulham & MF & Ball winning, low attacking involvement midfielders & 2.59 \\
Sandro Tonali & Newcastle & MF & Ball winning, low attacking involvement midfielders & 2.14 \\
Saša Lukić & Fulham & MF & Ball winning, low attacking involvement midfielders & 2.97 \\
Sean Longstaff & Newcastle & MF & Ball winning, low attacking involvement midfielders & 2.57 \\
Thomas Partey & Arsenal & MF & Ball winning, low attacking involvement midfielders & 2.90 \\
Tim Iroegbunam & Everton & MF & Ball winning, low attacking involvement midfielders & 4.93 \\
Tommy Doyle & Wolves & MF & Ball winning, low attacking involvement midfielders & 2.40 \\
Tomáš Souček & West Ham & MF & Ball winning, low attacking involvement midfielders & 3.81 \\
Tyler Adams & Bournemouth & MF & Ball winning, low attacking involvement midfielders & 3.23 \\
Tyler Dibling & Southampton & MF & Ball winning, low attacking involvement midfielders & 2.69 \\
Tyrick Mitchell & Crystal Palace & MF & Ball winning, low attacking involvement midfielders & 3.31 \\
Vitaly Janelt & Brentford & MF & Ball winning, low attacking involvement midfielders & 2.17 \\
Welington & Southampton & MF & Ball winning, low attacking involvement midfielders & 3.86 \\
Wes Burns & Ipswich Town & MF & Ball winning, low attacking involvement midfielders & 3.86 \\
Wilfred Ndidi & Leicester City & MF & Ball winning, low attacking involvement midfielders & 2.10 \\
Will Hughes & Crystal Palace & MF & Ball winning, low attacking involvement midfielders & 3.68 \\
William Smallbone & Southampton & MF & Ball winning, low attacking involvement midfielders & 3.41 \\
Yasin Ayari & Brighton & MF & Ball winning, low attacking involvement midfielders & 1.95 \\
Yehor Yarmoliuk & Brentford & MF & Ball winning, low attacking involvement midfielders & 2.16 \\
Yukinari Sugawara & Southampton & MF & Ball winning, low attacking involvement midfielders & 3.10 \\
Yves Bissouma & Tottenham & MF & Ball winning, low attacking involvement midfielders & 3.20 \\
Abdul Fatawu Issahaku & Leicester City & MF & High attacking involvement midfielders & 5.45 \\
Adama Traoré & Fulham & MF & High attacking involvement midfielders & 3.33 \\
Alejandro Garnacho & Manchester Utd & MF & High attacking involvement midfielders & 3.19 \\
Alex Iwobi & Fulham & MF & High attacking involvement midfielders & 2.34 \\
Alexis Mac Allister & Liverpool & MF & High attacking involvement midfielders & 4.79 \\
Amad Diallo & Manchester Utd & MF & High attacking involvement midfielders & 2.55 \\
Andreas Pereira & Fulham & MF & High attacking involvement midfielders & 4.18 \\
Anthony Elanga & Nottingham & MF & High attacking involvement midfielders & 4.70 \\
Antoine Semenyo & Bournemouth & MF & High attacking involvement midfielders & 3.40 \\
Bernardo Silva & Manchester City & MF & High attacking involvement midfielders & 2.96 \\
Brajan Gruda & Brighton & MF & High attacking involvement midfielders & 4.51 \\
Bruno Fernandes & Manchester Utd & MF & High attacking involvement midfielders & 3.45 \\
Bryan Mbeumo & Brentford & MF & High attacking involvement midfielders & 2.94 \\
Callum Hudson-Odoi & Nottingham & MF & High attacking involvement midfielders & 2.94 \\
Carlos Alcaraz & Everton / Southampton & MF & High attacking involvement midfielders & 4.16 \\
Christian Eriksen & Manchester Utd & MF & High attacking involvement midfielders & 5.66 \\
Christopher Nkunku & Chelsea & MF & High attacking involvement midfielders & 3.16 \\
Cody Gakpo & Liverpool & MF & High attacking involvement midfielders & 3.23 \\
Cole Palmer & Chelsea & MF & High attacking involvement midfielders & 3.63 \\
Crysencio Summerville & West Ham & MF & High attacking involvement midfielders & 3.52 \\
Curtis Jones & Liverpool & MF & High attacking involvement midfielders & 2.94 \\
Dango Ouattara & Bournemouth & MF & High attacking involvement midfielders & 3.50 \\
David Brooks & Bournemouth & MF & High attacking involvement midfielders & 3.02 \\
Declan Rice & Arsenal & MF & High attacking involvement midfielders & 3.32 \\
Dejan Kulusevski & Tottenham & MF & High attacking involvement midfielders & 2.19 \\
Diego Gómez & Brighton & MF & High attacking involvement midfielders & 5.42 \\
Dominik Szoboszlai & Liverpool & MF & High attacking involvement midfielders & 3.31 \\
Dwight McNeil & Everton & MF & High attacking involvement midfielders & 4.49 \\
Eberechi Eze & Crystal Palace & MF & High attacking involvement midfielders & 2.84 \\
Emile Smith Rowe & Fulham & MF & High attacking involvement midfielders & 3.02 \\
Enzo Fernández & Chelsea & MF & High attacking involvement midfielders & 2.23 \\
Facundo Buonanotte & Leicester City & MF & High attacking involvement midfielders & 3.85 \\
Georginio Rutter & Brighton & MF & High attacking involvement midfielders & 3.59 \\
Gonçalo Guedes & Wolves & MF & High attacking involvement midfielders & 2.90 \\
Harry Wilson & Fulham & MF & High attacking involvement midfielders & 3.33 \\
Iliman Ndiaye & Everton & MF & High attacking involvement midfielders & 4.21 \\
Ismaila Sarr & Crystal Palace & MF & High attacking involvement midfielders & 3.11 \\
Jack Grealish & Manchester City & MF & High attacking involvement midfielders & 4.82 \\
Jadon Sancho & Chelsea & MF & High attacking involvement midfielders & 3.50 \\
James Maddison & Tottenham & MF & High attacking involvement midfielders & 4.35 \\
Jarrod Bowen & West Ham & MF & High attacking involvement midfielders & 2.76 \\
Jeremy Doku & Manchester City & MF & High attacking involvement midfielders & 3.11 \\
Jota Silva & Nottingham & MF & High attacking involvement midfielders & 6.09 \\
Julio Enciso & Ipswich Town / Brighton & MF & High attacking involvement midfielders & 4.26 \\
Justin Devenny & Crystal Palace & MF & High attacking involvement midfielders & 4.23 \\
Justin Kluivert & Bournemouth & MF & High attacking involvement midfielders & 2.37 \\
Kaoru Mitoma & Brighton & MF & High attacking involvement midfielders & 2.46 \\
Kevin De Bruyne & Manchester City & MF & High attacking involvement midfielders & 5.25 \\
Kevin Schade & Brentford & MF & High attacking involvement midfielders & 3.76 \\
Leon Bailey & Aston Villa & MF & High attacking involvement midfielders & 3.04 \\
Marco Asensio & Aston Villa & MF & High attacking involvement midfielders & 2.77 \\
Marcus Rashford & Manchester Utd / Aston Villa & MF & High attacking involvement midfielders & 2.76 \\
Marcus Tavernier & Bournemouth & MF & High attacking involvement midfielders & 2.47 \\
Marshall Munetsi & Wolves & MF & High attacking involvement midfielders & 3.48 \\
Martin Ødegaard & Arsenal & MF & High attacking involvement midfielders & 3.56 \\
Mason Mount & Manchester Utd & MF & High attacking involvement midfielders & 4.72 \\
Matheus Cunha & Wolves & MF & High attacking involvement midfielders & 3.96 \\
Matt O'Riley & Brighton & MF & High attacking involvement midfielders & 2.71 \\
Mikel Merino & Arsenal & MF & High attacking involvement midfielders & 4.28 \\
Mikkel Damsgaard & Brentford & MF & High attacking involvement midfielders & 4.10 \\
Mohamed Salah & Liverpool & MF & High attacking involvement midfielders & 7.03 \\
Mohammed Kudus & West Ham & MF & High attacking involvement midfielders & 3.60 \\
Morgan Gibbs-White & Nottingham & MF & High attacking involvement midfielders & 2.01 \\
Morgan Rogers & Aston Villa & MF & High attacking involvement midfielders & 2.65 \\
Nathan Broadhead & Ipswich Town & MF & High attacking involvement midfielders & 3.77 \\
Noni Madueke & Chelsea & MF & High attacking involvement midfielders & 3.54 \\
Omar Marmoush & Manchester City & MF & High attacking involvement midfielders & 5.31 \\
Pablo Sarabia & Wolves & MF & High attacking involvement midfielders & 3.09 \\
Pedro Neto & Chelsea & MF & High attacking involvement midfielders & 2.78 \\
Phil Foden & Manchester City & MF & High attacking involvement midfielders & 2.92 \\
Reiss Nelson & Fulham / Arsenal & MF & High attacking involvement midfielders & 3.23 \\
Ross Barkley & Aston Villa & MF & High attacking involvement midfielders & 5.01 \\
Ryan Sessegnon & Fulham & MF & High attacking involvement midfielders & 6.16 \\
Sammie Szmodics & Ipswich Town & MF & High attacking involvement midfielders & 5.68 \\
Simon Adingra & Brighton & MF & High attacking involvement midfielders & 3.23 \\
Stephy Mavididi & Leicester City & MF & High attacking involvement midfielders & 2.81 \\
Sávio & Manchester City & MF & High attacking involvement midfielders & 3.64 \\
Tom Cairney & Fulham & MF & High attacking involvement midfielders & 5.50 \\
Yankuba Minteh & Brighton & MF & High attacking involvement midfielders & 2.81 \\
Youri Tielemans & Aston Villa & MF & High attacking involvement midfielders & 3.65 \\
İlkay Gündoğan & Manchester City & MF & High attacking involvement midfielders & 4.37 \\
\bottomrule
\end{longtable}


\section{Supplementary figures}
\label{sec:appendix:supplementary}

Figure~\ref{fig:corrheatmap} gives the full hierarchically clustered Spearman correlation
matrix over the outfield feature set, from which the divergence analysis in
Section~\ref{sec:results:divergence} is derived. Figure~\ref{fig:pcascree} reports the
scree and cumulative variance curves: \PCNinetyComponents{} components are required to
reach the ninety percent threshold, so the feature set is not close to low rank, and the
first three components together account for \PCCumThree{} of the variance.
Figure~\ref{fig:pcabiplot} shows the loadings biplot. Figure~\ref{fig:keepercorr} gives
the correlation structure of the surviving goalkeeper features and
Figure~\ref{fig:keepercompositedistributions} the distribution of the two goalkeeper
composites.

\begin{figure}[h]
  \centering
  \includegraphics[width=0.9\textwidth]{corr_heatmap_all.pdf}
  \caption{Hierarchically clustered Spearman correlation matrix over the outfield feature
  set in the earlier of the two reduced-sample seasons.}
  \label{fig:corrheatmap}
\end{figure}

\begin{figure}[h]
  \centering
  \begin{subfigure}[t]{0.48\textwidth}
    \includegraphics[width=\textwidth]{pca_scree.pdf}
    \caption{Scree and cumulative variance.}
    \label{fig:pcascree}
  \end{subfigure}
  \hfill
  \begin{subfigure}[t]{0.48\textwidth}
    \includegraphics[width=\textwidth]{pca_biplot.pdf}
    \caption{Loadings biplot.}
    \label{fig:pcabiplot}
  \end{subfigure}
  \caption{Principal component structure of the outfield feature set.}
  \label{fig:pcapanel}
\end{figure}

\begin{figure}[h]
  \centering
  \begin{subfigure}[t]{0.48\textwidth}
    \includegraphics[width=\textwidth]{keeper_corr_heatmap.pdf}
    \caption{Surviving goalkeeper features.}
    \label{fig:keepercorr}
  \end{subfigure}
  \hfill
  \begin{subfigure}[t]{0.48\textwidth}
    \includegraphics[width=\textwidth]{keeper_composite_distributions.pdf}
    \caption{Goalkeeper composite distributions.}
    \label{fig:keepercompositedistributions}
  \end{subfigure}
  \caption{Supplementary goalkeeper material, reduced sample.}
  \label{fig:keepersupplementary}
\end{figure}

\begin{figure}[h]
  \centering
  \begin{subfigure}[t]{0.48\textwidth}
    \includegraphics[width=\textwidth]{archive_keeper_corr.pdf}
    \caption{Correlation structure of the archive goalkeeper features.}
    \label{fig:arckeepercorr}
  \end{subfigure}
  \hfill
  \begin{subfigure}[t]{0.48\textwidth}
    \includegraphics[width=\textwidth]{archive_keeper_composites.pdf}
    \caption{Distribution of the archive goalkeeper composites.}
    \label{fig:arckeepercomposites}
  \end{subfigure}
  \caption{Supplementary goalkeeper material, archive sample. The full
  \KeeperFeatureCountArchive{} feature block supports the composite axes that the reduced
  sample cannot, which is why the archive keeper analysis in
  Section~\ref{sec:results:keepers} reports separations the reduced sample collapses.}
  \label{fig:arckeepersupplementary}
\end{figure}

\section{Individual archetype radars}
\label{sec:appendix:radars}

Figures~\ref{fig:arcradarsingledf} to \ref{fig:arcradarsinglefw} give the archive archetype
radars for each position group at full size, on the fixed axes described in
Section~\ref{sec:methods}, and Figures~\ref{fig:radarsingledf} to \ref{fig:radarsinglefw}
give the corresponding radars on the reduced sample. Within each position group the two
profiles are mirror images by construction, since the features are standardised within group
before clustering. The two sets are drawn on different feature counts, so their axes are not
directly comparable, and the comparison worth making is of shape rather than of extent.

\begin{figure}[h]
  \centering
  \includegraphics[width=0.68\textwidth]{archive_radar_DF.pdf}
  \caption{Defender archetypes, archive sample.}
  \label{fig:arcradarsingledf}
\end{figure}

\begin{figure}[h]
  \centering
  \includegraphics[width=0.68\textwidth]{archive_radar_MF.pdf}
  \caption{Midfielder archetypes, archive sample.}
  \label{fig:arcradarsinglemf}
\end{figure}

\begin{figure}[h]
  \centering
  \includegraphics[width=0.68\textwidth]{archive_radar_FW.pdf}
  \caption{Forward archetypes, archive sample.}
  \label{fig:arcradarsinglefw}
\end{figure}

\begin{figure}[h]
  \centering
  \includegraphics[width=0.68\textwidth]{radar_DF.pdf}
  \caption{Defender archetypes.}
  \label{fig:radarsingledf}
\end{figure}

\begin{figure}[h]
  \centering
  \includegraphics[width=0.68\textwidth]{radar_MF.pdf}
  \caption{Midfielder archetypes.}
  \label{fig:radarsinglemf}
\end{figure}

\begin{figure}[h]
  \centering
  \includegraphics[width=0.68\textwidth]{radar_FW.pdf}
  \caption{Forward archetypes.}
  \label{fig:radarsinglefw}
\end{figure}

\section{Further analyses on the reduced sample}
\label{sec:appendix:reducedanalyses}

The analyses collected here were run on the reduced sample of
Section~\ref{sec:data:reduced}. They are reported in full because several of them depend on
data that exists only for that window, and because the reduced sample is the robustness case
for the main result rather than a separate study. None of them carries the paper's argument,
which is why they sit here rather than in Section~\ref{sec:results}.

\subsection{Cluster selection and stability on the reduced sample}
\label{sec:results:clustering}

This subsection reports the clustering detail for the reduced sample of
Section~\ref{sec:data:reduced}, which is the robustness case rather than the main
analysis. Applying the same rule to the four selection criteria gives $k=\KGlobal{}$
globally and within each position group, matching the primary sample. The full criterion
curves appear in Figure~\ref{fig:kselection} and Appendix~\ref{sec:appendix:kselection}.

Stability separates the position groups sharply. Across \BootstrapN{} bootstrap resamples
the mean adjusted Rand index against the reference partition is \BootARIGlobal{} globally
and \BootARIMF{} for midfielders, both of which describe a partition that reappears almost
unchanged under resampling. For defenders the mean falls to \BootARIDF{} with a standard
deviation of \BootARISDDF{}, and for forwards to \BootARIFW{} with a standard deviation of
\BootARISDFW{}. The minimum adjusted Rand index observed across forward resamples is
\BootARIMinFW{}, which is to say that in the least favourable resample the partition
recovered is no better than chance. Forwards number only \NFWPrimary{} players above the
minutes threshold, so this is partly a sample size problem, but it is a real limit on what
can be claimed. The consensus co-assignment structure is shown in
Figure~\ref{fig:consensus}.

The relationship between discovered clusters and listed positions is weak. The global
partition scores an adjusted Rand index of \ARIvsPosition{} against the position
group with a cluster purity of \PurityVsPosition{}, so the clusters are not a rediscovery
of the positions that the source already records. Figure~\ref{fig:clustervsposition} shows
where the two disagree. Read together with the embedding geometry, the interpretation is
that the dominant axis in the surviving data is not defender against forward but
attacking involvement against defensive involvement, which cuts across the listed groups
and places a ball-playing full back nearer an ordinary midfielder than a defensive
counterpart at his own listed position.

\begin{figure}[t]
  \centering
  \includegraphics[width=\textwidth]{k_selection_curves.pdf}
  \caption{Cluster selection criteria across candidate cluster counts. The consensus rule
  was fixed before fitting and selects the same value in every scope.}
  \label{fig:kselection}
\end{figure}

\begin{figure}[t]
  \centering
  \begin{subfigure}[t]{0.48\textwidth}
    \includegraphics[width=\textwidth]{consensus_matrix.pdf}
    \caption{Bootstrap co-assignment.}
    \label{fig:consensus}
  \end{subfigure}
  \hfill
  \begin{subfigure}[t]{0.48\textwidth}
    \includegraphics[width=\textwidth]{cluster_vs_position.pdf}
    \caption{Clusters against listed positions.}
    \label{fig:clustervsposition}
  \end{subfigure}
  \caption{Cluster stability and the correspondence with listed positions.}
  \label{fig:clustervalidation}
\end{figure}

\begin{figure}[t]
  \centering
  \includegraphics[width=0.75\textwidth]{cluster_scatter_pca.pdf}
  \caption{The selected partition shown in the principal component space.}
  \label{fig:clusterscatter}
\end{figure}

\subsection{Descriptive structure}
\label{sec:results:descriptive}

After filtering, \NEligiblePrimary{} outfield players are eligible in the earlier of the two reduced-sample seasons,
comprising \NDFPrimary{} defenders, \NMFPrimary{} midfielders and \NFWPrimary{} forwards.
Summary statistics by position group appear in Table~\ref{tab:summary}, and the
distributions of four representative statistics are shown in
Figure~\ref{fig:distributions}.

The feature set is strongly non-normal. All \NFeaturesRejectNormal{} of the
\NFeatures{} features reject the Shapiro-Wilk null at the five percent level, mean
absolute skewness is \SkewMeanAbs{}, and the most skewed feature is goals per shot at
\SkewMax{}. Mardia's test rejects multivariate normality decisively on both the skewness
and kurtosis components. Principal component analysis requires no distributional
assumption to be computed, so this does not invalidate the decomposition, but it does
mean the components should be read as directions of variance rather than as latent
factors with any distributional interpretation, and it justifies the use of Spearman
rather than Pearson correlation throughout.

\begin{table}[t]
\centering
\small
\caption{Headline per-90 statistics by position group in the primary season, reported as mean with standard deviation in parentheses.}
\label{tab:summary}
\begin{tabular}{lrrr}
\toprule
Statistic & DF & MF & FW \\
\midrule
np\_xg\_p90 & 0.05 (0.04) & 0.16 (0.12) & 0.41 (0.17) \\
xa\_p90 & 0.06 (0.06) & 0.15 (0.10) & 0.16 (0.11) \\
key\_passes\_p90 & 0.51 (0.44) & 1.26 (0.63) & 1.05 (0.54) \\
xg\_chain\_p90 & 0.35 (0.16) & 0.50 (0.20) & 0.63 (0.20) \\
xg\_buildup\_p90 & 0.30 (0.14) & 0.28 (0.13) & 0.18 (0.09) \\
shots\_p90 & 0.54 (0.33) & 1.54 (0.85) & 2.53 (0.69) \\
shot\_accuracy\_pct & 29.22 (20.60) & 31.91 (13.77) & 40.30 (9.65) \\
crosses\_p90 & 1.22 (1.62) & 2.26 (1.75) & 1.40 (1.62) \\
interceptions\_p90 & 0.95 (0.37) & 0.74 (0.44) & 0.24 (0.16) \\
tackles\_won\_p90 & 1.06 (0.48) & 1.17 (0.55) & 0.55 (0.33) \\
fouls\_committed\_p90 & 0.87 (0.38) & 1.32 (0.55) & 1.27 (0.62) \\
offsides\_p90 & 0.06 (0.07) & 0.16 (0.17) & 0.53 (0.33) \\
\bottomrule
\end{tabular}
\end{table}


\begin{figure}[t]
  \centering
  \includegraphics[width=\textwidth]{distributions_by_position.pdf}
  \caption{Distributions of four representative per-90 statistics across position groups
  in the earlier of the two reduced-sample seasons. The separation between groups is clear for shot volume and for
  defensive actions, and much weaker for creation.}
  \label{fig:distributions}
\end{figure}

\subsection{Comparing embeddings}
\label{sec:results:embeddings}

Principal component analysis, UMAP and t-SNE were fitted to the same standardised
features and compared in Figure~\ref{fig:embeddings}. The selected UMAP configuration
used \UMAPBestNeighbors{} neighbours at a minimum distance of \UMAPBestMinDist{}, chosen
by trustworthiness over the twelve configuration grid reported in
Appendix~\ref{sec:appendix:umap}.

The three methods agree on the dominant structure and differ in how they present it.
Trustworthiness is \TrustPCA{} for the linear projection, \TrustUMAP{} for UMAP and
\TrustTSNE{} for t-SNE, while continuity is close to identical across the three. Nearest
neighbour position purity at ten neighbours is \PurityPCA{} for the linear projection and
\PurityUMAP{} for UMAP, so roughly two thirds to three quarters of a player's closest
statistical neighbours share his listed position group. That figure is informative in
both directions: it confirms that the surviving features carry real positional signal,
and it establishes that a substantial minority of every player's nearest neighbours are
listed at a different position, which is the space in which data-driven roles can differ
from listed ones.

All three embeddings display a continuous gradient from defenders through midfielders to
forwards rather than separated groups. That gradient, and not a set of islands, is the
honest description of the geometry, and it conditions how the clustering that follows
should be read.

\begin{figure}[t]
  \centering
  \includegraphics[width=\textwidth]{embedding_comparison.pdf}
  \caption{Two-dimensional embeddings of the same feature set, coloured by listed position
  group. The defender to forward gradient is continuous in all three, and t-SNE is shown
  for visual contrast only and is not used for clustering.}
  \label{fig:embeddings}
\end{figure}

\subsection{The third dimension}
\label{sec:results:threed}

Two components are a convenient way to draw the space and a poor way to describe it. The
first two carry shooting volume and wide creation and together account for
\ThreeDVarTwo{} of the variance. The third carries duels and discipline, loading on fouls
committed, cards and tackles won, and it is orthogonal to both of the others, so a planar
projection does not compress it, it deletes it. Retaining it raises captured variance to
\ThreeDVarThree{}. Figure~\ref{fig:threedpositions} shows the outfield space in all three,
drawn from two viewpoints so that structure can be told apart from the accident of a single
projection.

Adding the dimension does not change the headline geometry. The defender to forward gradient
remains continuous rather than resolving into islands, which is the same conclusion the
planar embeddings reached and is now established against a projection that had the
opportunity to falsify it. Figure~\ref{fig:threedclusters} places the selected partition in
the same space: the boundary cuts across the shooting axis and is close to indifferent to
duels and discipline, which is consistent with a split that is substantially about whether a
player shoots.

The third dimension does change the club picture, and this is the one place where it
produces a finding the planar analysis could not. Club centroids correlate with final
league position at \ThreeDTeamPCOne{} on the first component and \ThreeDTeamPCTwo{} on the
second, but only \ThreeDTeamPCThree{} on the third. Shooting volume and wide creation both
track league position; the physical dimension of duels, fouls and cards essentially does
not. A club does not finish higher by fouling more or less, which is unsurprising stated
plainly and is worth stating because the duelling axis is a full component of player
variation that carries almost no information about team success.
Figure~\ref{fig:threedteams} shows the clubs, labelled at the extremes of the table.

Finally, the three-dimensional view settles a structural question raised in
Section~\ref{sec:results:archetypes} by measurement rather than assertion. Drawing each
archetype centroid as a direction from the origin, the two centroids within a position
group have a cosine of \CosineDF{} for defenders, \CosineMF{} for midfielders and
\CosineFW{} for forwards. They are exactly anti-parallel, which is forced by standardising
within group before a two-way split. Figure~\ref{fig:threedaxes} shows the consequence: six
archetypes occupy three lines through the origin, not six independent directions.

\begin{figure}[t]
  \centering
  \includegraphics[width=\textwidth]{pca3d_positions.pdf}
  \caption{Outfield players in the first three principal components, from two viewpoints.
  Grey points on the floor are the vertical projection of each player, supplying a depth
  cue that a static rendering otherwise lacks.}
  \label{fig:threedpositions}
\end{figure}

\begin{figure}[t]
  \centering
  \includegraphics[width=\textwidth]{pca3d_clusters.pdf}
  \caption{The selected global partition in three components, with cluster centroids marked
  by crosses. The split runs along the shooting axis and is close to indifferent to the
  third component.}
  \label{fig:threedclusters}
\end{figure}

\begin{figure}[t]
  \centering
  \includegraphics[width=\textwidth]{pca3d_teams.pdf}
  \caption{Clubs as minutes-weighted centroids in three components, shaded by final league
  position with darker shading for a higher finish. Vertical lines to the floor give the
  depth cue, and only the top and bottom of the table are labelled.}
  \label{fig:threedteams}
\end{figure}

\begin{figure}[t]
  \centering
  \includegraphics[width=\textwidth]{pca3d_archetype_axes.pdf}
  \caption{Archetype centroid directions, drawn as unit vectors so the figure shows
  orientation rather than magnitude. Each position group forms a single line through the
  origin, so the six archetypes describe three axes.}
  \label{fig:threedaxes}
\end{figure}

\subsection{The same pipeline across the big five}
\label{sec:results:leagues}

Every result reported so far comes from one competition, which leaves an alternative
explanation standing: perhaps English football happens to divide in two and another league
would divide differently. The withdrawal applied to the whole site rather than to one
competition, so the other four of the big five are in the same reduced state and the
independent expected goals source covers all of them, which makes the alternative testable
rather than merely acknowledged.

Repeating the pipeline unchanged on Spain, Germany, Italy and France for the earlier of the two reduced-sample seasons
returns the same answer everywhere. The rule declared in advance selects
\LgKEng{} clusters in all \NLeagues{} leagues. The silhouette at that choice spans
\LgSilMin{} to \LgSilMax{} across the five, bootstrap stability spans \LgBootARIMin{} to
\LgBootARIMax{}, and agreement with listed position groups spans \LgARIPosMin{} to
\LgARIPosMax{}. Every one of those three ranges is narrow, and the third is uniformly weak:
in no league does the partition recover the positions the source already records.
Figure~\ref{fig:leagues} shows the comparison.

The elasticity result generalises as well. Fitted independently inside each league, all
\LgElastCount{} defensive elasticities lie between \LgElastMin{} and \LgElastMax{}, so
every one is below the value of one that the conventional possession adjustment assumes.
The overcorrection identified in Section~\ref{sec:results:padj} is therefore a property of
how defending responds to territory rather than a quirk of a single competition, and the
conventional correction overcorrects across the big five.

This changes the standing of the main result. Two clusters is not a fact about England. It
is what the surviving public feature set can support, and the same feature set supports the
same coarse division wherever it is applied.

\begin{figure}[t]
  \centering
  \includegraphics[width=\textwidth]{league_comparison.pdf}
  \caption{The pipeline applied unchanged to each of the big five leagues in the primary
  season. The selected number of clusters, the silhouette, bootstrap stability and
  agreement with listed positions are all close across competitions.}
  \label{fig:leagues}
\end{figure}

\subsection{What a supervised model sees}
\label{sec:results:supervised-reduced}

The clustering results leave one question open: are the surviving features too weak to
carry positional information at all, or do they carry it while the clustering chooses to
partition along a different axis? A supervised model settles it. Predicting the listed
position group from the same \NFeatures{} features, gradient boosting reaches an accuracy
of \LGBMAccuracy{} and a macro F1 of \LGBMMacroF{}, and logistic regression reaches
\LogitAccuracy{} and \LogitMacroF{}, against a majority class baseline of
\BaselineAccuracy{} and \BaselineMacroF{}. The features clearly do carry position, so the
weak agreement between clusters and positions reflects the clustering finding a different
and equally real axis, not an empty feature set. The confusion is almost entirely between
adjacent groups, as Figure~\ref{fig:confusion} shows, and forwards are the weakest class.

The SHAP decomposition in Figures~\ref{fig:shapbeeswarm} and \ref{fig:shapbyclass} carries
a finding that follows directly from the withdrawal. With possession and passing volume
gone, positions are identified largely by absence. Defenders are recognised by low
shooting and low key passes, forwards by low interceptions, and midfielders by an
intermediate profile on both. A classifier of this kind, trained before January 2026,
would have identified a defender by what he did with the ball. It now identifies him by
what he does not do.

Players the classifier places in the wrong group overlap significantly with the players
the clustering flagged as hybrids. The overlap rate is \OverlapLogit{} for logistic
regression at an enrichment of \EnrichmentLogit{} over independence, and
\OverlapLGBM{} at an enrichment of \EnrichmentLGBM{} for gradient boosting. The
enrichment is modest because the hybrid flag is itself broad, covering \NHybridPlayers{}
of \NEligiblePrimary{} players, so this is corroboration rather than a strong result.

The finishing above role regression illustrates a hazard worth recording. Predicting
non-penalty expected goals per ninety from non-shooting behaviour returns a cross
validated coefficient of determination of \FinishingRSq{}. An earlier specification that
also included xGChain returned \FinishingRSqLeaky{}, which looks far better and is in fact
circular: Understat credits xGChain for possessions the player finished himself, so the
difference between xGChain and xGBuildup contains the target and correlates with it at
\FinishingLeakCorr{}. The apparently strong model was reconstructing the outcome rather
than predicting it. Even at the honest value, roughly half the variance is unexplained, so
the residual leaderboard in Figure~\ref{fig:finishing} is presented as an illustration and
not as a measurement of finishing ability.

\begin{figure}[t]
  \centering
  \begin{subfigure}[t]{0.48\textwidth}
    \includegraphics[width=\textwidth]{confusion_matrix.pdf}
    \caption{Confusion between listed positions.}
    \label{fig:confusion}
  \end{subfigure}
  \hfill
  \begin{subfigure}[t]{0.48\textwidth}
    \includegraphics[width=\textwidth]{shap_by_class.pdf}
    \caption{Feature importance per position.}
    \label{fig:shapbyclass}
  \end{subfigure}
  \caption{Supervised validation of the surviving feature set.}
  \label{fig:supervised}
\end{figure}

\subsection{Team signatures}
\label{sec:results:teams}

Aggregating players into clubs produces the one result in this study that is both strong
and stable. The minutes-weighted position of a club's outfield players on the first
principal component correlates with final league position at a Spearman coefficient of
\TeamPCOneRho{} in the earlier of the two reduced-sample seasons and \TeamPCOneRhoRepl{} in the replication season,
which is close to an exact reproduction. Since the first component is essentially
attacking involvement, this says that clubs whose players collectively shoot and create
more finish higher, which is unsurprising in itself, but the near identical coefficient
across two independent seasons shows the team level summary is measuring something stable
even where the player level partitions are not.

Hierarchical clustering on archetype composition groups the \NTeams{} clubs into
\NStyleGroups{} style families, shown in Figure~\ref{fig:dendrogram}, with composition by
club in Figure~\ref{fig:composition}. Figure~\ref{fig:teampanels} places each club's
players in the shared embedding against the rest of the league, which makes squad shape
legible in a way that a composition bar cannot: clubs differ visibly in how far their
defenders sit toward the midfield region.

\begin{figure}[t]
  \centering
  \begin{subfigure}[t]{0.48\textwidth}
    \includegraphics[width=\textwidth]{team_dendrogram.pdf}
    \caption{Style families.}
    \label{fig:dendrogram}
  \end{subfigure}
  \hfill
  \begin{subfigure}[t]{0.48\textwidth}
    \includegraphics[width=\textwidth]{team_composition.pdf}
    \caption{Archetype composition by club.}
    \label{fig:composition}
  \end{subfigure}
  \caption{Club stylistic signatures, ordered by final league position.}
  \label{fig:teamsignatures}
\end{figure}

\begin{figure}[p]
  \centering
  \includegraphics[width=\textwidth]{team_panels.pdf}
  \caption{Each club's outfield players highlighted in the shared embedding, with the rest
  of the league in grey. Axes are shared across all twenty panels.}
  \label{fig:teampanels}
\end{figure}

\subsection{Similarity, replication and noise}
\label{sec:results:novel}

\subsubsection{Player similarity}
Nearest neighbour search within position group recovers football sense for players whose
roles are defined by the surviving statistics, and fails informatively for players whose
roles are not. Full backs return full backs, penalty box forwards return penalty box
forwards, and centre backs return centre backs. The instructive failures are the
midfielders. A holding midfielder's nearest neighbours are creative attacking midfielders,
because with no tackles by pitch third, no blocks, no clearances, no pressures and no
passing volume, the only things left that distinguish him are his set piece deliveries and
key passes. The metric finds the players who share that signature and misses the entire
defensive half of his job. An inverted full back returns three centre backs for the same
reason, since the role is defined by possession behaviour that no longer exists in the
data. Similarity scores among centre backs are systematically the lowest in the table,
which is the same feature poverty appearing as an inability to tell defenders apart at
all. Figure~\ref{fig:similarity} shows the structure.

\subsubsection{Cross-season replication}
\NCommonPlayers{} players are eligible in both seasons. The declared rule independently
selects the same number of clusters in every scope in both seasons, so the structure
replicates. Membership replicates only partly. The adjusted Rand index between seasons is
\CrossSeasonARIGlobal{} globally, with \CrossSeasonChangedGlobal{} players changing
cluster, and within position group it falls to \CrossSeasonARIDF{} for defenders,
\CrossSeasonARIMF{} for midfielders and \CrossSeasonARIFW{} for forwards. The within group
figures are worse than the bootstrap stability of Section~\ref{sec:results:clustering},
which is what would be expected if the within group splits are finer than the data can
reliably support. Figure~\ref{fig:crossseason} shows the transitions.

\subsubsection{Sensitivity to sampling noise}
\label{sec:results:shrinkage}
Empirical Bayes shrinkage toward the position group mean, with the degree of shrinkage set
by minutes played, provides the sharpest test of whether the partitions describe playing
roles or sampling noise. Shrinkage is substantial rather than cosmetic, and the results
divide cleanly.

The global partition survives. Only \ShrinkChangedGlobal{} of \NEligiblePrimary{} players
change cluster, the adjusted Rand index between the original and shrunken partitions is
\ShrinkARIGlobal{}, the selected number of clusters is unchanged, and the rate at which
players move falls monotonically with minutes played, exactly as the noise account
predicts.

The within group partitions fare worse, and they do not fare equally. For defenders the
adjusted Rand index falls to \ShrinkARIDF{}, with \ShrinkChangedDF{} of \NDFPrimary{}
players moving, which is a partition largely dissolved by removing sampling noise.
Midfielders retain \ShrinkARIMF{} and forwards \ShrinkARIFW{}, both of which survive in
part rather than collapsing. The forward figure is notably higher than it was before the
possession correction of Section~\ref{sec:results:padj} was applied, which is itself
informative: some of what previously looked like noise in the forward partition was
exposure to opponent possession, and removing it left a more reproducible split. The mechanism is
visible in the estimated variance components: several features are assigned an estimated
between-player variance of zero within a position group, meaning the model attributes
their entire observed spread to sampling noise and collapses them to the group mean,
removing them from the clustering altogether. Shot accuracy retains very little
reliability, because its true spread between players is small relative to the noise in
estimating it from a season of shots.

Read alongside the bootstrap and replication results, the three tests do not deliver a
single uniform verdict, and reporting them as though they did would overstate the case.
Defenders fail all three: bootstrap stability of \BootARIDF{}, cross-season agreement of
\CrossSeasonARIDF{}, and near dissolution under shrinkage. Forwards fail two of the three,
with a bootstrap minimum of \BootARIMinFW{} across resamples and cross-season agreement of
\CrossSeasonARIFW{}, while partly surviving shrinkage. Midfielders pass all three. The
conclusion that survives every test is therefore narrower than a single summary sentence
would suggest: the global division between attacking and defensive involvement is a real
property of players, the midfield split is trustworthy, and the defender split should not
be used to characterise an individual at all.
Figures~\ref{fig:shrinkage} and \ref{fig:shrinkageminutes} show the effect.

\begin{figure}[t]
  \centering
  \begin{subfigure}[t]{0.48\textwidth}
    \includegraphics[width=\textwidth]{similarity_heatmap.pdf}
    \caption{Nearest neighbour similarity.}
    \label{fig:similarity}
  \end{subfigure}
  \hfill
  \begin{subfigure}[t]{0.48\textwidth}
    \includegraphics[width=\textwidth]{cross_season_transitions.pdf}
    \caption{Between-season transitions.}
    \label{fig:crossseason}
  \end{subfigure}
  \caption{Similarity structure and cross-season replication.}
  \label{fig:novelpanel}
\end{figure}

\subsubsection{Role drift within a season}
\label{sec:results:drift}
Every result so far treats a season as one observation per player, which assumes a
player's profile is homogeneous across it. Match reports make that assumption testable,
and because the withdrawal removed the affected columns from the match schema rather than
emptying them, what remains there is fully populated. A full season is 380 reports at seven
seconds each, so the scrape is restricted to one club, \DriftTeam{}, chosen because its
eligible players span all six archetypes and therefore give drift the widest range to move
across. \DriftMatches{} matches were retrieved and \DriftPlayers{} players accumulated at
least one complete rolling window of \DriftWindow{} matches.

The match feature set has \DriftFeatures{} of the eighteen features used elsewhere, since
the Understat expected goals block has no match equivalent. Rather than assume that the
coarser space is a fair stand-in, we refit the season clustering using only those
\DriftFeatures{} features and compare it to the partition reported throughout. The two
agree at an adjusted Rand index of \DriftReducedARI{}, so the reduced space reproduces
almost all of the main division and drift measured within it is interpretable.

Most players do not drift. \DriftSwitchers{} of \DriftPlayers{}, or \DriftSwitchShare{}
percent, change cluster at least once during the season, and the remainder occupy a single
cluster throughout. The players who stay put are the ones whose roles are least ambiguous,
holding midfielders on one side and the recognised attackers on the other, while the movers
are predominantly full backs and wide players whose attacking involvement varies with the
fixture. Figure~\ref{fig:drift} shows the trajectories.

The finding cuts both ways and should be read that way. For four fifths of players a season
aggregate describes a single role and the homogeneity assumption holds. For the remaining
fifth it averages over two, which means their season-level archetype is a summary of
something that changed rather than a description of something stable. That is a limitation
of season aggregates in general rather than of this feature set in particular, and it is one
that only match level data can expose.

\begin{figure}[t]
  \centering
  \includegraphics[width=\textwidth]{role_drift.pdf}
  \caption{Rolling cluster assignment across the season for the \DriftTeam{} players who
  change cluster at least once, on a \DriftWindow{} match window. Players not shown occupied
  one cluster for the whole season.}
  \label{fig:drift}
\end{figure}

\subsubsection{Age and archetype}
\label{sec:results:age}
The two archetypes within a position group differ mainly in attacking involvement, and
involvement is plausibly something a career passes through rather than a fixed trait, so
the association with age is worth testing. It is tested within position group, because
archetypes are nested inside positions and a pooled table would largely recover the fact
that forwards and defenders have different age profiles.

Only defenders show a relationship. The crossing and chance creating defenders average
\AgeGapDF{} years younger than the low involvement group, a difference that the rank test
places at \AgeKruskalPDF{}, and the share of the crossing archetype falls monotonically
from the youngest band to the oldest. Midfielders show nothing at all, with a gap of
\AgeGapMF{} years at \AgeKruskalPMF{}, and the global partition is unrelated to age at
\AgeGlobalP{}. Forwards point the same way as defenders but cannot support the claim,
since six of the ten cells in their table have expected counts below five.

The defender result also illustrates a point about test choice that is easy to get wrong.
A chi-square on the banded table returns \AgeChiPDF{} and would ordinarily be read as no
association. That test spends its degrees of freedom on any departure from independence,
while the pattern here is monotonic, and a rank correlation against age recovers it at
\AgeTrendRhoDF{} with a p-value of \AgeTrendPDF{}. The banded test is not wrong, it is
underpowered against the alternative that actually holds, and reporting only the chi-square
would have converted a real finding into a null one. Figure~\ref{fig:age} shows the shares.

\begin{figure}[t]
  \centering
  \includegraphics[width=\textwidth]{age_archetype.pdf}
  \caption{Share of each age band taken by each archetype, within position group. The
  defender panel falls monotonically with age, which a chi-square on bands does not detect
  and a rank correlation does.}
  \label{fig:age}
\end{figure}

\begin{figure}[t]
  \centering
  \begin{subfigure}[t]{0.48\textwidth}
    \includegraphics[width=\textwidth]{shrinkage_effect.pdf}
    \caption{Effect of shrinkage on assignments.}
    \label{fig:shrinkage}
  \end{subfigure}
  \hfill
  \begin{subfigure}[t]{0.48\textwidth}
    \includegraphics[width=\textwidth]{shrinkage_by_minutes.pdf}
    \caption{Assignment changes by minutes played.}
    \label{fig:shrinkageminutes}
  \end{subfigure}
  \caption{Empirical Bayes shrinkage as a test of whether the partitions describe roles or
  noise.}
  \label{fig:shrinkagepanel}
\end{figure}

\begin{figure}[t]
  \centering
  \begin{subfigure}[t]{0.48\textwidth}
    \includegraphics[width=\textwidth]{outliers_scatter.pdf}
    \caption{Outliers within position group.}
    \label{fig:outliers}
  \end{subfigure}
  \hfill
  \begin{subfigure}[t]{0.48\textwidth}
    \includegraphics[width=\textwidth]{finishing_above_role.pdf}
    \caption{Residuals from the role model.}
    \label{fig:finishing}
  \end{subfigure}
  \caption{Unusual players, and the residual leaderboard that should not be read as a
  measurement of finishing ability.}
  \label{fig:outlierspanel}
\end{figure}

\begin{figure}[t]
  \centering
  \includegraphics[width=0.85\textwidth]{shap_beeswarm.pdf}
  \caption{Global SHAP values for the gradient boosted position classifier. Shooting
  volume and interceptions carry most of the signal, so a player's position is now
  identified largely by the actions he does not perform.}
  \label{fig:shapbeeswarm}
\end{figure}

\section{Reduced-sample supplementary figures}
\label{sec:appendix:reduced}

The figures collected here belong to the reduced sample of
Section~\ref{sec:data:reduced}. They are reported for completeness, because the reduced
sample is the robustness case for the main result rather than the main analysis itself.
Figure~\ref{fig:redarchetypeprofiles} gives the centroid deviations and
Figure~\ref{fig:redradarpanels} the archetype radars. Figures~\ref{fig:redkeeperclusters} to
\ref{fig:redkeeperranking} give the goalkeeper material that survives on that sample, which
is limited to shot stopping and workload for the reasons set out in
Section~\ref{sec:data:reduced}.

\begin{figure}[h]
  \centering
  \includegraphics[width=\textwidth]{archetype_profiles.pdf}
  \caption{Centroid deviation of each archetype from its position group mean, reduced
  sample.}
  \label{fig:redarchetypeprofiles}
\end{figure}

\begin{figure}[h]
  \centering
  \begin{subfigure}[t]{0.32\textwidth}
    \includegraphics[width=\textwidth]{radar_panel_DF.pdf}
    \caption{Defenders.}
  \end{subfigure}
  \hfill
  \begin{subfigure}[t]{0.32\textwidth}
    \includegraphics[width=\textwidth]{radar_panel_MF.pdf}
    \caption{Midfielders.}
  \end{subfigure}
  \hfill
  \begin{subfigure}[t]{0.32\textwidth}
    \includegraphics[width=\textwidth]{radar_panel_FW.pdf}
    \caption{Forwards.}
  \end{subfigure}
  \caption{Archetype radars on the reduced sample.}
  \label{fig:redradarpanels}
\end{figure}

\begin{figure}[h]
  \centering
  \begin{subfigure}[t]{0.48\textwidth}
    \includegraphics[width=\textwidth]{keeper_clusters_pca.pdf}
    \caption{Goalkeeper clusters in principal component space.}
    \label{fig:redkeeperclusters}
  \end{subfigure}
  \hfill
  \begin{subfigure}[t]{0.48\textwidth}
    \includegraphics[width=\textwidth]{keeper_composites.pdf}
    \caption{The two surviving composite axes.}
    \label{fig:redkeepercomposites}
  \end{subfigure}
  \caption{Goalkeeper structure on the reduced sample.}
  \label{fig:redkeeperstructure}
\end{figure}

\begin{figure}[h]
  \centering
  \begin{subfigure}[t]{0.48\textwidth}
    \includegraphics[width=\textwidth]{keeper_savepct_vs_workload.pdf}
    \caption{Save percentage against workload.}
    \label{fig:redkeeperworkload}
  \end{subfigure}
  \hfill
  \begin{subfigure}[t]{0.48\textwidth}
    \includegraphics[width=\textwidth]{keeper_saves_vs_sota.pdf}
    \caption{Saves against shots on target.}
    \label{fig:redkeepersota}
  \end{subfigure}
  \caption{Goalkeeper shot stopping on the reduced sample. With expected goals on target
  withdrawn, save percentage against volume is the only shot stopping view the reduced
  sample supports, and it cannot separate difficulty from quantity.}
  \label{fig:redkeepershotstopping}
\end{figure}

\begin{figure}[h]
  \centering
  \includegraphics[width=\textwidth]{keeper_shot_stopping_ranking.pdf}
  \caption{Goalkeeper shot stopping ranking on the reduced sample.}
  \label{fig:redkeeperranking}
\end{figure}
